%% ============================================================
%%  Chapter 5 — LLM Infrastructure: Reliability, Cost, and Configuration
%%  Production-Grade Agentic AI: LangGraph + Python Frameworks
%% ============================================================

\chapterquote{A system that talks to one model is a prototype. A system that
survives a model going down is a product.}{anon}

By the end of Chapter~4 your agent has a validated, strongly typed contract
between the model and the rest of the runtime. Every node that consumes model
output receives a Pydantic object, and every parse failure either repairs itself
or escalates cleanly. What the previous chapters deliberately left to one side
is the question of what happens \emph{below} that contract: who decides which
model to call, what happens when the call fails at the network or provider
level, and who pays attention to the bill.

These are not operational afterthoughts. In a production agent system the LLM
call is the single most expensive, slowest, and most failure-prone operation in
the hot path. Every model provider has a non-zero outage rate. Every model
version is eventually deprecated. Every run costs money, and budgets are finite.
If the infrastructure around that call is naive---a bare \code{litellm.completion()}
with no fallback and no budget guard---then the agent is one provider hiccup
away from a cascading failure and one runaway loop away from a costly bill.

This chapter builds that infrastructure in three layers. The first layer
\emph{routes} calls: a LiteLLM Router abstracts all provider calls behind a
single interface, pins exact model versions, and defines an ordered fallback
cascade. The second layer \emph{protects} calls: Tenacity decorators retry
transient faults per exception class, and a circuit breaker stops the runtime
from hammering a provider that is down for minutes rather than milliseconds.
The third layer \emph{governs} calls: a callback-driven cost tracker accumulates
real spend into \code{AgentState}, a budget circuit breaker converts an overage
into a clean graph stop, and a \code{FallbackRouter} logs every degradation
event to Langfuse. A final section externalises every model name, key, and
parameter into Pydantic Settings so that no hardcoded string appears in
application code.

The three custom exception classes introduced here---\code{BudgetExceededError},
\code{CircuitOpenError}, and \code{RetryExhaustedError}---map to three distinct
graph-edge outcomes. Table~\ref{tab:exception-routing} summarises that mapping
before you build any of the components, so you can keep the whole architecture
in view as each piece appears.

\begin{table}[ht]
\centering
\caption{Exception taxonomy and corresponding graph-edge outcomes}
\label{tab:exception-routing}
\begin{tabular}{lll}
\toprule
\textbf{Exception} & \textbf{Source} & \textbf{Graph edge} \\
\midrule
\code{BudgetExceededError} & \code{TokenBudget} validator & \code{budget\_exceeded} \\
\code{CircuitOpenError}    & \code{LLMCircuitBreaker}     & \code{circuit\_open}    \\
\code{RetryExhaustedError} & Tenacity \code{before\_sleep} & \code{retry\_exhausted} \\
\bottomrule
\end{tabular}
\end{table}

The sections of this chapter are ordered by dependency: first choose where
calls go (\S\ref{sec:litellm-router}), then decide which failures deserve
another attempt (\S\ref{sec:tenacity}), then decide when to stop calling a
broken provider entirely (\S\ref{sec:circuit-breaker}), then measure spend
(\S\ref{sec:cost-tracking}), then give overspend a clean graph exit
(\S\ref{sec:budget-circuit-breaker}), then make the cascade fully observable
(\S\ref{sec:fallback-routing}), and finally externalise every value that would
otherwise be a hardcoded string (\S\ref{sec:settings}). Each section hands a
resolved problem to the next.

The source layout for this chapter is:

\begin{verbatim}
src/
  config/
    llm_settings.py        # §5.7
    runnable_config.py     # §5.7
  llm/
    router_factory.py      # §5.1
    retry_policies.py      # §5.2
    circuit_breaker.py     # §5.3
    cost_tracker.py        # §5.4
    fallback_router.py     # §5.6
    exceptions.py          # §5.5
  graph/
    edges.py               # §5.5
    nodes.py               # §5.5
\end{verbatim}

%% ------------------------------------------------------------
\section{Why a Router Exists}
\label{sec:litellm-router}
%% ------------------------------------------------------------

The first question to resolve is why you need a router at all. A bare
\code{litellm.completion(model="gpt-4o", ...)} works perfectly well in
development, but it carries three production problems. First, it binds the call
site to a single provider. If that provider degrades, the call fails and there
is no recovery path. Second, the model string \code{"gpt-4o"} is an alias that
OpenAI can silently remap to a newer version at any time, breaking prompts that
were tuned against a specific checkpoint. Third, there is no single place to
configure timeouts, retries, or API base URLs; each call site accumulates its
own ad-hoc configuration.

LiteLLM's \code{Router} solves all three problems at once. You declare a
\code{model\_list}---a list of provider configurations each identified by a
logical alias---and a \code{fallbacks} dict that maps each alias to an ordered
list of fallback aliases. The router handles provider selection, request
dispatch, and failure routing internally. The call site uses only the logical
alias and never knows which provider actually fulfilled the request.

The model name in each entry's \code{litellm\_params.model} must be a fully
qualified, version-pinned string in the format \code{provider/model-version},
for example \code{openai/gpt-4o-2024-08-06} or
\code{anthropic/claude-3-5-sonnet-20241022}. Using a bare alias like
\code{gpt-4o} is the same mistake as not pinning a dependency version: it works
until the day the provider remaps the alias, and debugging a prompt regression
caused by a silent model upgrade is expensive. Version pinning makes that class
of failure impossible.

\begin{lstlisting}[language=Python,
  caption={RouterFactory assembling a LiteLLM Router from externalised settings},
  label={lst:router-factory}]
# src/llm/router_factory.py

from __future__ import annotations
from typing import Any
from litellm import Router
from src.config.llm_settings import LLMSettings


def _make_model_entry(alias: str, model: str, api_key: str) -> dict[str, Any]:
    return {
        "model_name": alias,
        "litellm_params": {
            "model":   model,
            "api_key": api_key,
        },
    }


class RouterFactory:
    """Assemble a LiteLLM Router entirely from LLMSettings.

    No model name, API key, or routing parameter is hardcoded here.
    All values are injected from settings, which are loaded from the
    environment (see §5.7).
    """

    @staticmethod
    def build(settings: LLMSettings) -> Router:
        import litellm
        litellm.set_verbose = settings.debug

        model_list = [
            _make_model_entry(
                alias="primary",
                model=settings.primary_model,
                api_key=settings.openai_api_key.get_secret_value(),
            ),
            _make_model_entry(
                alias="fallback_1",
                model=settings.fallback_model_1,
                api_key=settings.openai_api_key.get_secret_value(),
            ),
            _make_model_entry(
                alias="fallback_2",
                model=settings.fallback_model_2,
                api_key=settings.anthropic_api_key.get_secret_value(),
            ),
        ]

        # Priority routing: requests always attempt aliases in the order
        # they appear in model_list before considering the fallbacks dict.
        return Router(
            model_list=model_list,
            fallbacks=[{"primary": ["fallback_1", "fallback_2"]}],
            routing_strategy="priority",
            num_retries=0,        % retries are handled by §5.2 decorators
            timeout=settings.request_timeout_sec,
        )
\end{lstlisting}

Two decisions in Listing~\ref{lst:router-factory} deserve explanation.
\code{num\_retries=0} disables LiteLLM's built-in retry so that Section~5.2's
Tenacity decorators remain the single source of retry behaviour. Duplicate retry
logic at two layers compounds backoff timing and makes failure analysis
confusing. The \code{routing\_strategy="priority"} choice is deliberate: for
cost-sensitive agents you want predictable ordering---primary first, cheaper
fallback second---rather than latency-optimised load balancing. Chapter~8
revisits this choice for high-throughput deployments where \code{"least-busy"}
is more appropriate.

\begin{warningbox}
Never pass \code{api\_key} as a plain string literal anywhere in your source
tree. The pattern \code{settings.openai\_api\_key.get\_secret\_value()} is the
only correct unwrapping site. Pydantic's \code{SecretStr} field prevents the key
from appearing in logs, \code{repr()}, or serialised state. Section~5.7 explains
the full settings hierarchy.
\end{warningbox}

The router centralises \emph{where} calls go. The next question is
\emph{which failures deserve another attempt} --- and the answer depends
entirely on the type of failure, not on a single global retry count.

%% ------------------------------------------------------------
\section{Retries Depend on Failure Type}
\label{sec:tenacity}
%% ------------------------------------------------------------

Retry logic is not a single policy. A \code{RateLimitError} is retriable
because the provider will accept the request again once the rate window resets;
waiting and retrying is the correct response. An \code{AuthenticationError} is
not retriable because sending the same invalid credentials again will produce the
same failure; retrying wastes budget and delays the error signal. A catch-all
\code{except Exception} wrapper silently conflates the two: it retries auth
failures, swallows the error, and burns tokens doing it.

The correct pattern is one \code{@retry} decorator per exception class, each
with parameters calibrated to that class's failure semantics. Tenacity supports
this directly through \code{retry\_if\_exception\_type}, which composes cleanly
with \code{wait\_exponential} and \code{stop\_after\_attempt}.

\begin{lstlisting}[language=Python,
  caption={Per-exception Tenacity decorators with calibrated backoff parameters},
  label={lst:retry-policies}]
# src/llm/retry_policies.py

from __future__ import annotations
import logging
from tenacity import (
    retry,
    wait_exponential,
    stop_after_attempt,
    retry_if_exception_type,
    before_sleep_log,
    RetryError,
)
from litellm.exceptions import RateLimitError, APIConnectionError, Timeout
from litellm import Router, ModelResponse
from src.llm.exceptions import BudgetExceededError, RetryExhaustedError

logger = logging.getLogger(__name__)


# --- RateLimitError ---
# Five attempts: provider rate windows typically reset within 60 s.
# Exponential backoff from 2 s to 60 s covers both short and long windows.
@retry(
    retry=retry_if_exception_type(RateLimitError),
    wait=wait_exponential(multiplier=1, min=2, max=60),
    stop=stop_after_attempt(5),
    before_sleep=before_sleep_log(logger, logging.WARNING),
    reraise=False,
)
def _call_rate_limited(router: Router, payload: dict) -> ModelResponse:
    return router.completion(**payload)


# --- APIConnectionError ---
# Three attempts: connection failures are often transient (DNS blip,
# TCP reset) but if three consecutive attempts fail, the provider is
# likely down and the circuit breaker (§5.3) should open instead.
@retry(
    retry=retry_if_exception_type(APIConnectionError),
    wait=wait_exponential(multiplier=1, min=1, max=30),
    stop=stop_after_attempt(3),
    before_sleep=before_sleep_log(logger, logging.WARNING),
    reraise=False,
)
def _call_connection_retry(router: Router, payload: dict) -> ModelResponse:
    return router.completion(**payload)


# --- Timeout ---
# Two attempts only: a timeout already consumed wall-clock time equal to
# the configured request_timeout_sec. Two attempts means at most 2x that
# cost before giving up.
@retry(
    retry=retry_if_exception_type(Timeout),
    wait=wait_exponential(multiplier=1, min=5, max=30),
    stop=stop_after_attempt(2),
    before_sleep=before_sleep_log(logger, logging.WARNING),
    reraise=False,
)
def _call_timeout_retry(router: Router, payload: dict) -> ModelResponse:
    return router.completion(**payload)


def resilient_call(router: Router, payload: dict) -> ModelResponse:
    """Dispatch a completion through the appropriate retry wrapper.

    BudgetExceededError is deliberately not wrapped in any @retry
    decorator: it must propagate immediately to the graph's error-routing
    edge (§5.5). Retrying a budget-exceeded call would only deepen the
    overage.
    """
    if isinstance(payload.get("__exc_hint"), RateLimitError):
        try:
            return _call_rate_limited(router, payload)
        except RetryError as exc:
            raise RetryExhaustedError("RateLimitError") from exc

    try:
        return _call_connection_retry(router, payload)
    except RetryError as exc:
        raise RetryExhaustedError("APIConnectionError") from exc
\end{lstlisting}

The \code{before\_sleep\_log} callback writes a structured log entry on every
sleep, including the attempt number, exception type, and calculated wait
duration. These entries feed the cost-tracking callback in Section~5.4 via the
same structured logging pipeline: each retry is a signal that budget may be
burning without progress.

\begin{notebox}
\code{reraise=False} means Tenacity raises a \code{RetryError} wrapping the last
exception when all attempts are exhausted, rather than re-raising the original
exception class directly. The \code{resilient\_call} wrapper converts
\code{RetryError} to the domain-specific \code{RetryExhaustedError} so the graph
edge routing in Section~5.5 can distinguish it from a first-attempt failure.
\end{notebox}

Retries handle faults that resolve within seconds. The next failure mode is
different in kind: a provider that is unhealthy for minutes. Retrying against
that provider burns the full Tenacity budget on every call. The answer is a
circuit breaker that fails immediately once the provider is known to be down,
and probes cautiously for recovery.

%% ------------------------------------------------------------
\section{When to Stop Calling a Broken Provider}
\label{sec:circuit-breaker}
%% ------------------------------------------------------------

Tenacity retries handle transient failures. A circuit breaker handles
\emph{sustained} failures---a provider that is down for minutes, not
milliseconds. Without a circuit breaker, retrying against a down provider burns
the full Tenacity budget on every call, adding seconds of latency to every
agent cycle. With a circuit breaker, the first burst of failures opens the
circuit, subsequent calls fail immediately without touching the network, and the
circuit probes for recovery after a timeout.

There are Python libraries that implement circuit breakers, but using one here
would hide the state machine behind an API surface. The implementation below is
short enough to read in five minutes and long enough to contain all the
production-relevant behaviour: thread-safe state transitions, a half-open probe
that allows exactly one test call, and an extended recovery timeout on repeated
failures.

The three states are:

\begin{itemize}
  \item \textbf{CLOSED} --- the normal operating state. Calls pass through.
        Failures are counted. When the failure count reaches
        \code{failure\_threshold}, the circuit transitions to OPEN.
  \item \textbf{OPEN} --- the tripped state. Calls fail immediately with
        \code{CircuitOpenError} without touching the provider. After
        \code{recovery\_timeout} seconds, the circuit transitions to HALF\_OPEN.
  \item \textbf{HALF\_OPEN} --- the probing state. Exactly one call is allowed
        through. If it succeeds, the circuit closes and the failure count resets.
        If it fails, the circuit reopens with a doubled recovery timeout.
\end{itemize}

\begin{lstlisting}[language=Python,
  caption={Circuit breaker as an explicit three-state machine},
  label={lst:circuit-breaker}]
# src/llm/circuit_breaker.py

from __future__ import annotations
import threading
import time
from enum import Enum, auto
from typing import Any, Callable
from litellm import ModelResponse
from src.llm.exceptions import CircuitOpenError


class CircuitState(Enum):
    CLOSED    = auto()
    OPEN      = auto()
    HALF_OPEN = auto()


class LLMCircuitBreaker:
    """
    Thread-safe circuit breaker for a single LiteLLM router tier.

    State machine (ASCII diagram for inline reference):
      CLOSED  --[failures >= threshold]--> OPEN
      OPEN    --[recovery_timeout elapsed]--> HALF_OPEN
      HALF_OPEN --[probe succeeds]--> CLOSED
      HALF_OPEN --[probe fails]----> OPEN (timeout *= 2)
    """

    def __init__(
        self,
        failure_threshold: int = 5,
        recovery_timeout: float = 60.0,
    ) -> None:
        self.failure_threshold  = failure_threshold
        self._base_timeout      = recovery_timeout
        self._current_timeout   = recovery_timeout
        self.state              = CircuitState.CLOSED
        self._failure_count     = 0
        self._last_failure_time: float | None = None
        self._lock              = threading.Lock()

    def call(
        self,
        fn: Callable[..., ModelResponse],
        *args: Any,
        **kwargs: Any,
    ) -> ModelResponse:
        with self._lock:
            if self.state == CircuitState.OPEN:
                if self._should_attempt_reset():
                    self._transition(CircuitState.HALF_OPEN)
                else:
                    raise CircuitOpenError(
                        f"Circuit is OPEN; recovery in "
                        f"{self._seconds_until_reset():.1f}s"
                    )

        try:
            result = fn(*args, **kwargs)
            self._record_success()
            return result
        except Exception as exc:
            self._record_failure()
            raise

    def _transition(self, new_state: CircuitState) -> None:
        self.state = new_state

    def _should_attempt_reset(self) -> bool:
        if self._last_failure_time is None:
            return False
        elapsed = time.monotonic() - self._last_failure_time
        return elapsed >= self._current_timeout

    def _seconds_until_reset(self) -> float:
        if self._last_failure_time is None:
            return 0.0
        elapsed = time.monotonic() - self._last_failure_time
        return max(0.0, self._current_timeout - elapsed)

    def _record_success(self) -> None:
        with self._lock:
            if self.state == CircuitState.HALF_OPEN:
                self._failure_count   = 0
                self._current_timeout = self._base_timeout
                self._transition(CircuitState.CLOSED)
            elif self.state == CircuitState.CLOSED:
                self._failure_count = 0

    def _record_failure(self) -> None:
        with self._lock:
            self._failure_count    += 1
            self._last_failure_time = time.monotonic()

            if self.state == CircuitState.HALF_OPEN:
                self._current_timeout *= 2
                self._transition(CircuitState.OPEN)
            elif (
                self.state == CircuitState.CLOSED
                and self._failure_count >= self.failure_threshold
            ):
                self._transition(CircuitState.OPEN)
\end{lstlisting}

Two implementation choices are worth explaining in detail.

First, \code{time.monotonic()} rather than \code{datetime.now()} is used for
timing. Monotonic clocks are guaranteed never to go backwards, which matters
because system clocks can jump during NTP corrections or when a container is
migrated. A backwards jump on a wall clock would cause the circuit to remain
open far longer than \code{recovery\_timeout}. Monotonic time has no such
vulnerability.

Second, the actual call to \code{fn} happens \emph{outside} the lock. If the
lock were held during the I/O call, every thread in the process would block
waiting for the lock to release, defeating the purpose of multithreading
entirely. The state check and transition happen under the lock; the network call
does not.

Chapter~11 revisits this class to produce an \code{async}-compatible variant
using \code{asyncio.Lock} and \code{asyncio.get\_event\_loop().run\_in\_executor}
for high-concurrency deployments.

With routing, retries, and circuit breaking in place, the runtime now has
defences against provider faults. The remaining operational risk is financial:
the agent could run for many cycles and accumulate a large bill before any node
notices. The next section addresses that by making cost visible in real time,
directly inside \code{AgentState}.

%% ------------------------------------------------------------
\section{Tracking Spend in Real Time}
\label{sec:cost-tracking}
%% ------------------------------------------------------------

The natural place to accumulate cost is at the call site --- a line like
\code{budget += estimate\_tokens(prompt) * price\_per\_token} after each
\code{router.completion()} call. That approach has a fatal flaw: you do not know
the actual token count until the response arrives, per-model pricing is not
a constant you should hardcode, and the estimate can be wrong in ways that
compound silently across cycles.

LiteLLM already solves this. It emits a \code{success\_callback} after every
completed call, passes the full response object including token usage, and
exposes \code{litellm.completion\_cost()} which knows the current price schedule
for each supported model. This is the correct accumulation point: it fires
exactly once per successful call, after the response has been validated, with
accurate usage data.

The \code{CostTracker} class registers a method as that callback and accumulates
cost into a \code{TokenBudget} Pydantic model. \code{TokenBudget} lives inside
\code{AgentState} so that every node in the graph can read the current spend
without importing the tracker directly. When an accumulation would push
\code{used\_usd} past \code{limit\_usd}, the \code{TokenBudget} validator raises
\code{BudgetExceededError} immediately, before the update is committed.

\begin{lstlisting}[language=Python,
  caption={TokenBudget model and CostTracker callback wiring},
  label={lst:cost-tracker}]
# src/llm/cost_tracker.py

from __future__ import annotations
import threading
import litellm
from pydantic import BaseModel, model_validator
from litellm import ModelResponse
from src.llm.exceptions import BudgetExceededError


class TokenBudget(BaseModel):
    used_usd:                float = 0.0
    limit_usd:               float
    prompt_tokens_total:     int   = 0
    completion_tokens_total: int   = 0

    @model_validator(mode="after")
    def check_budget(self) -> "TokenBudget":
        if self.used_usd > self.limit_usd:
            raise BudgetExceededError(
                used_usd=self.used_usd,
                limit_usd=self.limit_usd,
            )
        return self


class CostTracker:
    """Registers itself as a LiteLLM success_callback and accumulates cost.

    Registration is per-Router instance, not global, to avoid cross-agent
    pollution in multi-agent deployments where multiple Routers coexist in
    the same process.
    """

    def __init__(self, budget: TokenBudget, router: litellm.Router) -> None:
        self._budget = budget
        self._lock   = threading.Lock()
        router.success_callback = [self.success_callback]

    def success_callback(
        self,
        kwargs:              dict,
        completion_response: ModelResponse,
        start_time:          float,
        end_time:            float,
    ) -> None:
        call_cost = litellm.completion_cost(
            completion_response=completion_response
        )
        usage = completion_response.usage
        with self._lock:
            self._budget = TokenBudget(
                used_usd=self._budget.used_usd + call_cost,
                limit_usd=self._budget.limit_usd,
                prompt_tokens_total=(
                    self._budget.prompt_tokens_total
                    + (usage.prompt_tokens if usage else 0)
                ),
                completion_tokens_total=(
                    self._budget.completion_tokens_total
                    + (usage.completion_tokens if usage else 0)
                ),
            )

    def get_budget_snapshot(self) -> TokenBudget:
        with self._lock:
            return self._budget.model_copy()
\end{lstlisting}

The \code{AgentState} type now carries two new fields to support cost tracking
and the downstream sections of this chapter:

\begin{lstlisting}[language=Python,
  caption={New AgentState fields for cost tracking and degradation},
  label={lst:agent-state-ch5}]
# src/agent/state.py (additions)

from src.llm.cost_tracker import TokenBudget

class AgentState(TypedDict):
    # ... existing fields from previous chapters ...
    token_budget:       TokenBudget        # updated by CostTracker
    last_model_used:    str                # populated by FallbackRouter (§5.6)
    last_call_cost_usd: float              # scalar for Langfuse event
    degraded:           bool               # True when static fallback used
    stop_reason:        StopReason | None  # None while running
\end{lstlisting}

\begin{notebox}
The \code{TokenBudget} validator raises \code{BudgetExceededError} at
construction time, which means the exception surfaces inside
\code{success\_callback}, which runs inside LiteLLM's response handling. The
\code{CostTracker} therefore catches \code{BudgetExceededError}, stores it on
an internal flag, and lets the graph node check \code{get\_budget\_snapshot()}
immediately after a call completes. Section~5.5 shows the graph-side wiring that
reads this flag and routes accordingly.
\end{notebox}

Now that spend is visible inside every graph cycle, the runtime needs a policy
for what an overage means at the graph level. The answer is not a retry and not
an exception that bubbles up through the runner. It is a clean, structured
graph stop---the subject of the next section.

%% ------------------------------------------------------------
\section{Treating Budget Overage as a Graph Exit}
\label{sec:budget-circuit-breaker}
%% ------------------------------------------------------------

A budget overage is not a network error or a model failure. It is a business
rule violation: the agent has consumed more than the operator authorised. The
correct response is not a retry, not a fallback to a cheaper model, and not an
unhandled exception that propagates up through the graph runner. It is a clean,
observable stop that writes the final state, records the reason, and returns a
structured signal to the caller.

Three components implement this: a custom exception, a graph routing function,
and a terminal node.

\begin{lstlisting}[language=Python,
  caption={BudgetExceededError and the StopReason model},
  label={lst:budget-exception}]
# src/llm/exceptions.py

from __future__ import annotations
from typing import Literal
from pydantic import BaseModel


class BudgetExceededError(Exception):
    def __init__(self, used_usd: float, limit_usd: float) -> None:
        super().__init__(
            f"Budget exceeded: used={used_usd:.4f} USD, "
            f"limit={limit_usd:.4f} USD"
        )
        self.used_usd  = used_usd
        self.limit_usd = limit_usd


class CircuitOpenError(Exception):
    pass


class RetryExhaustedError(Exception):
    def __init__(self, exception_class: str) -> None:
        super().__init__(
            f"All retries exhausted for exception class: {exception_class}"
        )
        self.exception_class = exception_class


class StopReason(BaseModel):
    code:     Literal["budget_exceeded", "circuit_open", "max_iterations"]
    detail:   str
    used_usd: float | None = None
\end{lstlisting}

\begin{lstlisting}[language=Python,
  caption={Graph edge and terminal node for budget-exceeded routing},
  label={lst:budget-graph}]
# src/graph/edges.py

from __future__ import annotations
from typing import Literal
from src.agent.state import AgentState


def should_continue_or_stop(
    state: AgentState,
) -> Literal["continue", "budget_exceeded"]:
    """Conditional edge after every LLM call node.

    Reads the budget snapshot from state rather than calling the tracker
    directly, so that the graph remains a pure function of state.
    """
    budget = state.get("token_budget")
    if budget is None:
        return "continue"
    if budget.used_usd >= budget.limit_usd:
        return "budget_exceeded"
    return "continue"


# src/graph/nodes.py (budget handler)

from langfuse import Langfuse
from langgraph.types import RunnableConfig
from src.agent.state import AgentState
from src.llm.exceptions import StopReason


def handle_budget_exceeded(
    state: AgentState,
    config: RunnableConfig,
) -> AgentState:
    """Terminal node: flush state, emit Langfuse event, return stop reason.

    BudgetExceededError is not re-raised here. It has already been caught
    and logged by the CostTracker. This node's job is to produce a clean
    final state.
    """
    langfuse = Langfuse()
    budget   = state.get("token_budget")
    used     = budget.used_usd if budget else 0.0
    limit    = budget.limit_usd if budget else 0.0

    langfuse.event(
        name="budget_exceeded",
        metadata={
            "used_usd":  used,
            "limit_usd": limit,
            "model":     state.get("last_model_used", "unknown"),
        },
    )

    return {
        **state,
        "stop_reason": StopReason(
            code="budget_exceeded",
            detail=(
                f"Agent stopped: used {used:.4f} USD of "
                f"{limit:.4f} USD budget."
            ),
            used_usd=used,
        ),
    }
\end{lstlisting}

The graph wiring connects these two pieces with a conditional edge that runs
after every node that performs an LLM call:

\begin{lstlisting}[language=Python,
  caption={Graph wiring for the budget-exceeded branch},
  label={lst:budget-wiring}]
# src/agent/graph.py (excerpt)

from langgraph.graph import END, StateGraph
from src.graph.edges import should_continue_or_stop
from src.graph.nodes import handle_budget_exceeded

b.add_node("handle_budget_exceeded", handle_budget_exceeded)

b.add_conditional_edges(
    "reason_node",
    should_continue_or_stop,
    {
        "continue":        "validation_node",
        "budget_exceeded": "handle_budget_exceeded",
    },
)

b.add_edge("handle_budget_exceeded", END)
\end{lstlisting}

The \code{StopReason} object in \code{AgentState} is the authoritative signal
for any caller of the compiled graph. A calling orchestrator, API handler, or
test checks \code{final\_state["stop\_reason"]} rather than catching an
exception, because the graph guarantees it always terminates with a structured
final state rather than propagating an error out of the runner. Chapter~6
persists \code{stop\_reason} to episodic memory so that retrospective analysis
can distinguish budget stops from normal completions. Chapter~12 surfaces it in
the operator dashboard.

The budget exit works, but it operates on a single tier. The production system
need a wider defensive layer: one that tries cheaper providers before giving up,
and that logs every tier attempt for post-mortem analysis. That is the job of
the fallback cascade.

%% ------------------------------------------------------------
\section{Fallback Cascade with Full Observability}
\label{sec:fallback-routing}
%% ------------------------------------------------------------

LiteLLM's built-in fallback mechanism is convenient for simple cases, but it
operates at the library level and is invisible to your application's
observability stack. When the primary model fails and LiteLLM silently routes to
the fallback, you have no event in Langfuse, no update to \code{AgentState}, and
no opportunity to check whether the budget permits the fallback call. The
\code{FallbackRouter} class makes the cascade explicit: each tier is a named
\code{FallbackTier} object with its own circuit breaker, and every
attempt---successful or not---is logged as a Langfuse event with a consistent
schema.

\begin{lstlisting}[language=Python,
  caption={FallbackRouter with per-tier circuit breakers and Langfuse logging},
  label={lst:fallback-router}]
# src/llm/fallback_router.py

from __future__ import annotations
from dataclasses import dataclass
from typing import Union
import logging
from langfuse import Langfuse
from litellm import ModelResponse
from langgraph.types import RunnableConfig
from pydantic import BaseModel
from src.llm.circuit_breaker import LLMCircuitBreaker
from src.llm.cost_tracker import CostTracker
from src.llm.exceptions import CircuitOpenError, RetryExhaustedError
from src.llm.retry_policies import resilient_call
from src.agent.state import AgentState

logger = logging.getLogger(__name__)


@dataclass
class FallbackTier:
    model_alias:      str
    circuit_breaker:  LLMCircuitBreaker
    cost_ceiling_usd: float


class StaticFallbackResponse(BaseModel):
    content:    str  = "Service temporarily unavailable. Please retry shortly."
    degraded:   bool = True
    model_used: str  = "static"
    cost_usd:   float = 0.0


class FallbackRouter:
    """Explicit three-tier fallback cascade with per-attempt Langfuse events.

    Each tier wraps its own LLMCircuitBreaker instance so that failures on
    the primary tier do not contaminate the fallback tier's health counter.
    """

    def __init__(
        self,
        tiers:    list[FallbackTier],
        tracker:  CostTracker,
        langfuse: Langfuse,
        router,
    ) -> None:
        self._tiers    = tiers
        self._tracker  = tracker
        self._langfuse = langfuse
        self._router   = router

    def route(
        self,
        payload: dict,
        state:   AgentState,
        config:  RunnableConfig,
    ) -> Union[ModelResponse, StaticFallbackResponse]:
        budget    = self._tracker.get_budget_snapshot()
        remaining = budget.limit_usd - budget.used_usd

        for tier_index, tier in enumerate(self._tiers):
            if remaining < tier.cost_ceiling_usd:
                self._log_attempt(
                    tier=tier, tier_index=tier_index,
                    success=False, cost_usd=0.0,
                    reason="InsufficientBudget",
                )
                continue

            call_payload = {**payload, "model": tier.model_alias}
            try:
                response: ModelResponse = tier.circuit_breaker.call(
                    resilient_call, self._router, call_payload
                )
                call_cost = float(
                    getattr(response, "_hidden_params", {}).get(
                        "response_cost", 0.0
                    )
                )
                self._log_attempt(
                    tier=tier, tier_index=tier_index,
                    success=True, cost_usd=call_cost,
                    reason=None,
                )
                state["last_model_used"]    = tier.model_alias
                state["last_call_cost_usd"] = call_cost
                return response

            except (CircuitOpenError, RetryExhaustedError) as exc:
                reason = type(exc).__name__
                logger.warning(
                    "Tier %d (%s) failed: %s",
                    tier_index, tier.model_alias, reason,
                )
                self._log_attempt(
                    tier=tier, tier_index=tier_index,
                    success=False, cost_usd=0.0,
                    reason=reason,
                )
                continue

        self._log_attempt(
            tier=FallbackTier("static", None, 0.0),  # type: ignore[arg-type]
            tier_index=len(self._tiers),
            success=False, cost_usd=0.0,
            reason="AllTiersExhausted",
        )
        state["degraded"] = True
        return StaticFallbackResponse()

    def _log_attempt(
        self,
        tier:       FallbackTier,
        tier_index: int,
        success:    bool,
        cost_usd:   float,
        reason:     str | None,
    ) -> None:
        self._langfuse.event(
            name="llm_call_attempt",
            metadata={
                "model_used": tier.model_alias,
                "tier":       tier_index,
                "success":    success,
                "cost_usd":   cost_usd,
                "reason":     reason,
            },
        )
\end{lstlisting}

The \code{reason} field in each Langfuse event is always the exception class
name, never free text. This keeps Langfuse queries filterable: you can build a
dashboard query for all events where \code{reason == "CircuitOpenError"} and
immediately see how often the primary tier is tripping the breaker. Free-text
reasons would make that aggregation impossible. Chapter~7 covers the full
Langfuse trace hierarchy and shows how \code{llm\_call\_attempt} events nest
under the parent agent trace. Chapter~12 covers the alerting rules that fire
when consecutive \code{success=False} events exceed a threshold.

With all three infrastructure layers in place---routing, protection, and
governance---the last risk is configuration drift: model names, API keys, and
budget limits that live in source code rather than the environment. The next
section eliminates that risk entirely.

%% ------------------------------------------------------------
\section{Configuration Must Not Live in Code}
\label{sec:settings}
%% ------------------------------------------------------------

Every section in this chapter has deferred one question: where do the model
names, API keys, budget limits, and retry parameters come from? The answer is a
single \code{LLMSettings} class that inherits from Pydantic's
\code{BaseSettings} and loads all values from environment variables prefixed
with \code{LLM\_}. Nothing else in the source tree knows a model name as a
string literal.

This matters for three reasons. First, the same codebase can run against
different model configurations in development, staging, and production by
changing environment variables rather than source code. Second, secrets never
appear in code, logs, or version control because \code{SecretStr} fields prevent
them from being serialised or included in \code{repr()}. Third, a CI pipeline
can validate the production \code{.env} file against the \code{LLMSettings}
schema before deployment, catching a missing or malformed variable before it
causes a runtime error. Chapter~13 shows that CI validation step in detail.

\begin{lstlisting}[language=Python,
  caption={LLMSettings loading all configuration from the environment},
  label={lst:llm-settings}]
# src/config/llm_settings.py

from __future__ import annotations
import re
from functools import lru_cache
from pydantic import SecretStr, field_validator
from pydantic_settings import BaseSettings, SettingsConfigDict


_MODEL_PIN_RE = re.compile(r"^[a-z0-9_\-]+/[a-z0-9\-\.]+$")


class LLMSettings(BaseSettings):
    """All LiteLLM router configuration, loaded from environment variables.

    Variable names are the field names uppercased and prefixed with LLM_.
    Example .env.development:

        LLM_PRIMARY_MODEL=openai/gpt-4o-2024-08-06
        LLM_FALLBACK_MODEL_1=openai/gpt-4o-mini-2024-07-18
        LLM_FALLBACK_MODEL_2=anthropic/claude-3-5-haiku-20241022
        LLM_OPENAI_API_KEY=sk-...
        LLM_ANTHROPIC_API_KEY=sk-ant-...
        LLM_TOTAL_BUDGET_USD=2.00
        LLM_FAILURE_THRESHOLD=5
        LLM_RECOVERY_TIMEOUT_SEC=60.0
        LLM_REQUEST_TIMEOUT_SEC=30.0
        LLM_DEBUG=false
    """

    model_config = SettingsConfigDict(
        env_prefix="LLM_",
        env_file=(".env", ".env.production"),
        secrets_dir="/run/secrets",
        case_sensitive=False,
    )

    primary_model:          str
    fallback_model_1:       str
    fallback_model_2:       str
    openai_api_key:         SecretStr
    anthropic_api_key:      SecretStr
    total_budget_usd:       float
    failure_threshold:      int   = 5
    recovery_timeout_sec:   float = 60.0
    request_timeout_sec:    float = 30.0
    debug:                  bool  = False

    @field_validator(
        "primary_model", "fallback_model_1", "fallback_model_2",
        mode="before",
    )
    @classmethod
    def validate_model_pin(cls, v: str) -> str:
        """Reject bare aliases like 'gpt-4o'; require 'openai/gpt-4o-2024-08-06'."""
        if not _MODEL_PIN_RE.match(v):
            raise ValueError(
                f"Model name must be in 'provider/model-version' format, "
                f"got: {v!r}. Using bare aliases risks silent version drift."
            )
        return v


@lru_cache(maxsize=1)
def get_settings() -> LLMSettings:
    """Return the process-scoped LLMSettings singleton.

    The lru_cache ensures the .env file is read exactly once per process,
    which matters for performance and for test isolation (tests can
    invalidate the cache with get_settings.cache_clear()).
    """
    return LLMSettings()
\end{lstlisting}

\begin{lstlisting}[language=Python,
  caption={Building a RunnableConfig with per-invocation overrides},
  label={lst:runnable-config}]
# src/config/runnable_config.py

from __future__ import annotations
from typing import Any
from langgraph.types import RunnableConfig
from src.config.llm_settings import LLMSettings, get_settings


def build_runnable_config(
    settings:  LLMSettings | None = None,
    overrides: dict[str, Any] | None = None,
) -> RunnableConfig:
    """Produce a RunnableConfig that carries LLM settings into graph nodes.

    The 'configurable' dict is the LangGraph-standard location for
    per-invocation parameters. Nodes access it via:
        config["configurable"].get("model_override")

    Overrides allow a single graph invocation to use a different model
    tier than the default without mutating the global settings. This is
    used by the test suite (§9.1) to force specific tiers.
    """
    resolved = settings or get_settings()
    configurable: dict[str, Any] = {
        "primary_model":        resolved.primary_model,
        "fallback_model_1":     resolved.fallback_model_1,
        "fallback_model_2":     resolved.fallback_model_2,
        "total_budget_usd":     resolved.total_budget_usd,
        "failure_threshold":    resolved.failure_threshold,
        "recovery_timeout_sec": resolved.recovery_timeout_sec,
        "request_timeout_sec":  resolved.request_timeout_sec,
    }
    if overrides:
        configurable.update(overrides)
    return RunnableConfig(configurable=configurable)
\end{lstlisting}

The \code{LLM\_} prefix namespace prevents collision with other environment
variables in the same deployment. In a multi-service container environment where
a dozen services share an environment, an unprefixed variable named
\code{PRIMARY\_MODEL} is a liability. Every settings class in this system uses
its own prefix; Chapter~8 adds an \code{EMBED\_} prefix for embedding model
settings and a \code{RERANK\_} prefix for reranker settings, following the same
pattern.

The dependency direction through this chapter is now fully linear and explicit:
\code{get\_settings()} provides \code{LLMSettings} to \code{RouterFactory.build()},
which produces a \code{Router} for \code{CostTracker} and \code{FallbackRouter}.
\code{FallbackRouter} is injected into graph nodes. No component reaches across
the chain; each layer only knows its direct dependency. Chapter~9 exploits this
by constructing the chain in test fixtures with synthetic settings and a mock
router, making every component testable in isolation without network access.
